\section{Related work}\label{sec:related}

\subsection{Trace clustering}\label{subsec:rw-clustering}

Mining a single model from a heterogeneous log yields a structure that describes no actual
behaviour, and the established remedy is to partition the log first and discover a specialised
model for each group \citep{greco2006clustering}. Partitioning needs a representation. Trace
profiles over activities, transitions, originators and case attributes make vector-space
clustering applicable to event data, and different profiles yield materially different partitions
\citep{song2009trace}; features built from conserved subsequences restore the order information
that bag-of-activities profiles discard \citep{bose2009context}; learned dense representations of
activities and traces outperform hand-crafted profiles on the same task
\citep{dekoninck2018act2vec}. A systematic review organises fifteen years of this work into four
interchangeable dimensions, namely representation, algorithm, use of non-control-flow information,
and evaluation strategy, and it already lists downstream model quality among the recognised
evaluation strategies \citep{zandkarimi2020generic}. The clearest instance is ActiTraC, which
admits a trace into a cluster only when the fitness of the process model discovered for that
cluster does not significantly degrade, so that the downstream artefact drives the grouping rather
than a distance in feature space \citep{deweerdt2013actitrac}. What that literature calls a
variant is close to what we call a process type \citep{taymouri2021variant}.

The difference here is one of placement. ActiTraC uses the downstream score inside the clustering
loop, as a wrapper; we use it afterwards, on partitions that four algorithm families produced
independently over two text arms, which is what allows the question of whether the algorithm
matters at all to be posed. We evaluate emergent process-type partitions by the predictive
performance of per-cluster sequence models, extending the model-driven evaluation principle of
\citet{deweerdt2013actitrac} from process-model fitness to next-activity prediction, and adding
shuffled-partition and global-model controls that this literature does not routinely report.

\subsection{Per-cluster sequence models}\label{subsec:rw-percluster}

Learning one sequence model per cluster is not a recent arrangement. \citet{ferreira2007sequence}
treat clustering as a mixture of first-order Markov chains, so that each cluster \emph{is} a
sequence model and a trace belongs to the cluster whose model most probably generated it;
clustering and per-cluster estimation are a single expectation--maximisation procedure.
\citet{difrancescomarino2019clusteringppm} separate the two stages: prefixes of completed cases
are clustered by control flow, a classifier is trained for each cluster, and a running case is
routed at runtime to the classifier of the cluster it matches. How the sequence is encoded is a
first-order determinant of how such per-cluster predictors perform \citep{teinemaa2019outcome},
and a cross-field comparison of next-element prediction methods finds that the techniques aiming
at interpretable models are the ones that give up accuracy \citep{tax2020interdisciplinary}.

We claim none of this. Our departures are narrow, and each is stated as such: third-order Markov
chains alongside a directly-follows graph and three neural families, rather than a single
first-order chain; complete concluded cases rather than prefixes; a separate extraction step that
turns each per-cluster model into one inspectable path; and, in place of the assumption that the
arrangement pays for itself, a shuffled-partition null and a global-model baseline that measure
how much of the apparent gain survives. Prediction is a measuring instrument here, not the
product.

\subsection{Representative and canonical traces}\label{subsec:rw-representative}

Summarising a cluster by a single execution has formal machinery behind it. Multiple-sequence
alignment, borrowed from bioinformatics, aligns traces column-wise so that the conserved core and
the deviations can be read off directly \citep{bose2012diagnostics}. \citet{chatain2017alignment}
go further and make model traces the centroids of the clusters, so that model explanations of
observed behaviour drive the partition instead of vector-space similarity; later work generalises
the centroid from a single run to a sub-model \citep{boltenhagen2021modelbased}. The idea of one
canonical execution per cluster belongs to that line. A single path is also a deliberate move to
the simplicity corner of the fitness/precision/generalisation/simplicity frame, bought at the cost
of generalisation \citep{buijs2014quality}, and the directly-follows graph we use as one of the
path learners carries its own documented distortions \citep{vanderaalst2019dfg}.

What differs is what the method is allowed to assume. Alignment-based centroids require a process
model as input, and the absence of any documented as-is model is the premise of this study rather
than an inconvenience within it. Our representative is derived from a learned sequence model over
the cluster's own traces, carries a likelihood, is comparable across two trace alphabets and five
model families, and is scored against the medoid of the cluster's real traces, which we report as
the strong baseline it turns out to be. The alignment consensus remains the alternative
representative we did not compute, and we say so where it bears on the interpretation of our
results.

\subsection{Predictive process monitoring and its critiques}\label{subsec:rw-ppm}

Next-activity prediction from event logs is mature. A single recurrent network predicts the next
activity, its timestamp and the remaining case suffix \citep{tax2017lstm}; treating traces as
language licenses the sequence-modelling apparatus wholesale \citep{evermann2017predicting};
attention-based encoders followed \citep{bukhsh2021processtransformer}. The benchmarks matter more
than the architectures. Under a common protocol, reported superiority frequently fails to
reproduce and no architecture dominates across logs \citep{ramamaneiro2023dlppm}, and the splits
in standard use leak information between training and test in ways that inflate reported
performance \citep{weytjens2022leakage}. Both observations are reasons to read a single accuracy
figure as a statement about a particular log and protocol rather than about a method, and they set
the standard of evidence we hold our own comparison to.

The critique we take most seriously is that of \citet{peeperkorn2023rnnstructure}, who show that a
network reaching good next-activity accuracy can still fail to capture the structure of the
underlying process, particularly when generalising to unseen variants. That result shapes our
design rather than threatening it. Fitness and next-activity objectives can disagree, so one
cannot be substituted for the other without argument. We therefore use accuracy retrospectively,
as a measure of how learnable a cluster is, report path-support measures alongside it, and draw no
conclusion about the process itself from an accuracy figure.

\subsection{Cluster validity}\label{subsec:rw-validity}

The indices we compute were introduced as \emph{relative} instruments, that is, devices for
ranking candidate partitions produced under comparable conditions, not absolute measures of
quality \citep{rousseeuw1987silhouettes, daviesbouldin1979measure, calinski1974dendrite,
moulavi2014dbcv}. Read instead as objectives, they misbehave: partitions that directly maximise
popular indices frequently disagree with expert reference labellings \citep{gagolewski2021invalid}.
Over 2.7 million partitions, \citet{yerbury2025rvi} find relative validity indices empirically
unsuitable for choosing between clustering approaches, and recommend outcome-oriented criteria
informed by domain requirements. A further asymmetry bears directly on cross-algorithm
comparison: no established index evaluates whether noise assignments are \emph{correct}, so a
density-based clustering that emits noise and a partitional one that labels every point are not
scored on the same population \citep{beer2026disco}. The genre is familiar from topic modelling,
where a field's consensus intrinsic metric proved negatively correlated with human judgement of
the same models \citep{chang2009teahaves}, and from the argument that clustering admits no
application-independent quality procedure at all \citep{vonluxburg2012clustering}.

The counter-case is \citet{chicco2025dbcv}, who find DBCV more informative than three competing
density-based indices on synthetic data and on health records. We do not contest that finding. It
compares indices within one paradigm on recovery of cluster geometry, where agreement with an
external index reached 60\%, and informativeness about geometry is a different quantity from
predictiveness of downstream utility. Our contribution to this thread is an instance rather than a
theory: sixteen partitions of one administrative corpus, with the indices and a downstream measure
reported side by side, and the resulting null stated at the strength sixteen partitions can
support.

\subsection{Privacy and public-sector process mining}\label{subsec:rw-privacy-public}

Event logs are personal data in their own right, because a trace is a behavioural fingerprint that
survives the removal of names \citep{mannhardt2019privacy, elkoumy2022privacy}. The response has
been sanitisation designed around what the downstream analysis needs: $k$-anonymity over trace
variants \citep{fahrenkrogpetersen2019pretsa}, extension to the contextual attributes attached to
events \citep{fahrenkrogpetersen2020pripel}, and an optimisation formulation of the
privacy--utility frontier \citep{fahrenkrogpetersen2023optimal}. Closest to our own experiment,
\citet{ferrao2022utility} re-cluster an anonymised administrative cohort and compare the resulting
solutions with three of the indices we use. Prior work therefore quantifies the utility cost of
event-log sanitisation for process discovery \citep{fahrenkrogpetersen2019pretsa,
fahrenkrogpetersen2023optimal}, of anonymisation for clustering low-dimensional tabular microdata
\citep{ferrao2022utility}, and of text de-identification for supervised NLP
\citep{lothritz2023evaluating}; to our knowledge the effect of free-text de-identification on
unsupervised trace clustering quality and on next-activity prediction has not been reported, and
it is that narrow question we answer against a pre-registered margin.

Process mining has been applied to public administration and to courts, and to auditing before
either \citep{vanderaalst2010auditing, jans2014fieldstudy}. Two precedents sit closest.
\citet{unger2021jurimetrics} build an event log from business-law lawsuits at a Brazilian state
court and find that proceedings follow an ad hoc rather than a designed sequence flow;
\citet{nai2026procurement} reconstruct public-procurement workflows from official European tender
records, extract case attributes from legal texts with language models, and have the discovered
models validated by domain experts. Brazilian work has likewise argued for process mining as a
transparency instrument and analysed an electronic process system in the judiciary
\citep{araujo2023mineracao, limaesilva2026analise}, alongside e-government case studies and
adoption studies elsewhere \citep{delgado2022discovery, racis2025donot, trottier2025using}. Every
one of these analyses a procedure whose identity is known in advance. Our unit of analysis is the
set of procedure types itself, because the administration's own declared taxonomy does not record
it reliably, and the catalogue that results is the object we then put under evaluation.

\subsection{Transparency, auditability and administrative procedure}\label{subsec:rw-legal}

Scholarship on administrative transparency has largely given up the equation of disclosure with
accountability. \citet{ananny2018seeing} set out why seeing a system amounts neither to knowing
nor to governing it, and \citet{janssen2012benefits} find the same confusion in open government,
where publication is routinely mistaken for use. \citet{bovens2007analysing} supplies the
vocabulary the distinction needs: accountability is a relation in which an actor explains conduct
to a forum competent to question and to judge. A catalogue on that reading is not a disclosure
but the explanatory content such a forum requires, which makes naming the forum part of the
claim rather than an afterthought.

This venue's own literature treats administrative procedure at that level.
\citet{suksi2021administrative} argues, from one national legal order, that due-process
safeguards written for human civil servants thin out as procedure is automated, and that the
safeguards worth having are preventive rather than reactive, which is an argument for making an
as-is procedure inspectable before anything is automated on top of it.
\citet{branting2025sensitive} work on government records under freedom-of-information
obligations with a system in operational use at federal agencies, and \citet{tinarrage2025empirical}
assess an institutional device of Brazilian law empirically, pairing classification results with a
legal reading of why it underperforms. Closest in method and jurisdiction at once, \citet{araujo2025trace} apply trace clustering to
Brazilian judicial process data in this journal, for the purpose of simplifying those
procedures and identifying patterns in them. We were unable to obtain the full text and do not
compare results against it; we record it because it is the nearest published pairing of this
technique with this jurisdiction, and a reader assessing novelty should start there. The
difference in setting is nonetheless substantive: judicial procedure is externally codified,
so a partition of it can be checked against a procedural code, whereas the administrative
routing studied here has no such external specification --- which is precisely why the record
has to be derived rather than looked up.
Closest to our method on the evidence we could read, \citet{vianna2024topic} discover
topic structure without supervision over Brazilian Portuguese legal documents and have the
result assessed by legal specialists, a step our catalogue has not been through
(Section~\ref{sec:threats}). Two constraints follow. What the records contain decides which task
may be attempted on them \citep{medvedeva2023rethinking}, and a taxonomy offered for legal use
is more defensible when produced by a stated procedure with explicit ending conditions
\citep{nickerson2013taxonomy}.
